\addbibresource{references.bib}

\title{Voting in Elder-led Churches: Six Duties}
\author{Jason Reeves (June 2026)}
\date{}

\begin{document}

\ifdefined\embedversion
  \pagestyle{empty}
\fi

\maketitle
\extrafloats{100}

\begin{abstract}
This paper is a synthesis of teaching from Alexander Strauch in his book, \textit{Recovering Biblical Eldership}\footcite{strauch2000recovering} as well as material from 9Marks.  The paper is from the perspective of churches which substantially adhere to the Second London Baptist Confession of Faith (1689) but has implications for all Christian churches working to exhibit biblical eldership in an age in which such a thing is more and more rare and unpopular.  What are the duties of church members and elders with regard to voting on church matters?  When and how is it permitted to vote against your elders?  This paper will attempt to explore these questions.
\newline
\newline
The six duties outlined in this paper with regard to voting on church matters are: \textit{Duty to Defer}, \textit{Duty of Diligence}, \textit{Duty to Declare}, \textit{Duty to Decline}, \textit{Duty to Depart}, and the \textit{Duty to Disencumber}.
\end{abstract}

\newthought{There is an old 9Marks article} from 2014 with a humorous title: \textit{I Move We Don't Vote So Much!}\footcite{gilbert2014}  Baptists in general seem to have a real affinity for voting.  I am thankful that voting on minutiae has not been a staple of the church I've served for nearly a decade.  In Thom Rainer's \textit{Autopsy of A Deceased Church}, published in the same year as the 9Marks article, he focuses on "minutiae meetings" as one of the marks of a church in decline.\footcite{rainer2014}  But even in churches like mine where we don't engage in minutiae, people can be very touchy about voting - \textit{what's} being voted on, \textit{how} we go about voting on it, etc.  I believe many times people get the wrong impression about voting on church matters -- that it is similar to voting in a government election, that they can just "make up their mind" at the last minute and "let their conscience be their guide" even without having done any due diligence to find out how God would have them vote on a matter.

Much less do many consider that there are certain \textit{duties} that Christians have and must carry out with regard to how they would cast their votes.  In a world rife with the spirit of rugged individualism and self absorption, many believe that their vote is merely some "sacred" expression of their own desires.\footnote{Note that this is not necessarily the same thing as "voting your conscience".  In order for a "conscience vote" to be meaningful, your conscience must be educated with the correct information and sanctified by the Holy Spirit.  Otherwise, "voting your conscience" simply means "voting your feelings" or "do what you like".}  To vote against their true desires almost seems to them like an assault on their own personhood.

In today's world, \textit{duty}, the idea that we humans have obligations to other human beings and ultimately to God, seems like an antiquated concept.  However, God assures us that, as with all matters of His law, it is an eternal concept.\footnote{The end of the matter; all has been heard. Fear God and keep his commandments, \textbf{for this is the whole duty of man}. For God will bring every deed into judgment, with every secret thing, whether good or evil. \\ \hspace*{\fill} --Ecclesiastes 12:13-14 (\textsc{esv})} In this paper, I will attempt to describe the duties Christian church members have with regard to voting on church matters, basing my arguments on Scripture and referencing other materials when necessary.

You may notice that in this paper I dwell heavily on the topic of church officer votes, and especially voting for the office of elder.  This is more of a "feature" than a "bug".  The most important and crucial matters that come before church members for a vote are the ones dealing with who will lead.  Next in importance are the votes dealing with the "keys of the kingdom"\footnote{Matthew 16:18-19} - voting to allow people into the church's membership, and voting to dismiss members, whether to another body or as the result of a church discipline proceeding.  Budgetary and lesser matters, while they may be more numerous in Baptist business meetings, and may actually require a vote, are not really the focus of this paper.  The duties we'll discuss \textit{do} apply to these lesser matters, but these matters pale in comparison to officer elections.  Let's begin with \textit{Duty to Defer}.

\ifdefined\embedversion\newpage\fi
\section{Duty to Defer}
Christian church members have a \textit{duty to defer} to their elders' leadership.  This doesn't mean that you must defer if your conscience is troubled or if you are clearly being asked to do something unscriptural.  Rather, if the disagreement is simply a matter of not seeing \textit{exactly}\footnote{Minor disagreements should be possible to set aside.  Major ones require further action and communication.  See \textit{Duty of Diligence} and \textit{Duty to Declare}} eye-to-eye, the duty of the Christian church member is to vote how the elders have requested.  Paul says to us in I Thessalonians:
\begin{quote}
But we request of you, brethren, that you appreciate those who diligently labor among you, and have charge over you in the Lord and give you instruction, and that you esteem them very highly in love because of their work. Live in peace with one another.\footnote{I Thessalonians 5:12-13 \textsc{(nasb 77)}}
\end{quote}

\noindent
The writer of Hebrews makes the point a bit more sharp:
\begin{quote}
Obey your leaders, and submit to them; for they keep watch over your souls, as those who will give an account. Let them do this with joy and not with grief, for this would be unprofitable for you.\footnote{Hebrews 13:17 \textsc{(nasb 77)}}
\end{quote}

\noindent
Obviously, the work of appreciating, esteeming very highly in love, obeying, and submitting to elders does not consist of following elders blindly.  Both J.C. Ryle and Charles Spurgeon have used the old adage "the best of men are men at best", although the origin is with neither of them.  However, it is no less true despite its questionable heritage.  The best men among us are but poor sinners, and will tell you so.  They have blind spots.  They will fail.\footnote{Present company included.}

Also, I do not mean to diminish the role and responsibility of the congregation to examine, for example, men put forth as church officers.  Both Acts 6 and I Timothy 3 seem to indicate that congregational examination is an important part of choosing men for office.

The 1689 Confession in Chapter 21, Paragraph 2, says:
\begin{quote}
God alone is Lord of the Conscience\footnote{James 4:12; Romans 14:4}, and hath left it free from the Doctrines and Commandments of men,\footnote{Acts 4:19, 5:29; I Corinthians 7:23; Matthew 15:9} which are in any thing contrary to his Word, or not contained in it. So that to Believe such Doctrines, or obey such Commands out of Conscience,\footnote{Colossians 2:20,22,23} is to betray true liberty of Conscience; and the requiring of an implicit Faith,\footnote{I Corinthians 3:5; II Corinthians 1:24} and absolute and blind Obedience, is to destroy Liberty of Conscience, and Reason also.\footnote{Scripture references from the Confession itself.}
\end{quote}

\noindent
So we have liberty of conscience, and of course we must use our liberty wisely.  Paul tells us in Galatians 5:
\begin{quote}
For you were called to freedom, brethren; only \textbf{do not turn your freedom into an opportunity for the flesh, but through love serve one another}. For the whole Law is fulfilled in one word, in the statement, "YOU SHALL LOVE YOUR NEIGHBOR AS YOURSELF." But if you bite and devour one another, take care lest you be consumed by one another.\footnote{Galatians 5:13-15 \textsc{(nasb 77)}}
\end{quote}

\noindent
And to the Ephesians:
\begin{quote}
And do not get drunk with wine, for that is dissipation, but be filled with the Spirit, speaking to one another in psalms and hymns and spiritual songs, singing and making melody with your heart to the Lord; always giving thanks for all things in the name of our Lord Jesus Christ to God, even the Father; and \textbf{be subject to one another in the fear of Christ}.\footnote{Ephesians 5:18-21 \textsc{(nasb 77)}} 
\end{quote}

\noindent
So you see we, as Christians, have \textit{true liberty} in Christ.  But we are warned not to use our liberty in order to satisfy the desires of our flesh, but rather to serve one another in love.  To further complicate things is this principle of "mutual submission" propounded in Ephesians.  It seems like a contradiction.  How can two Christians \textit{both} be submitted to one another?  One of them has to win, right?  

I would point out that in a marriage, this contradiction is resolved because God has made the husband the head.  He is the tie breaker.  When one brother wants option A and another wants option B, who wins?  We are taught to consider the needs of others more highly than our own needs.\footnote{If therefore there is any encouragement in Christ, if there is any consolation of love, if there is any fellowship of the Spirit, if any affection and compassion, make my joy complete by being of the same mind, maintaining the same love, united in spirit, intent on one purpose. Do nothing from selfishness or empty conceit, but with humility of mind \textbf{let each of you regard one another as more important than himself; do not merely look out for your own personal interests, but also for the interests of others}.  \\ \hspace*{\fill} --Philippians 2:1-4 (\textsc{nasb 77})}  But if two brothers cannot agree on who is going to defer to the other, God has ordained that the elders are to lead.  When the sheep are scattered and uncertain, it's up to the elders to bring them together and to lead them, \textit{but} it's also up to the sheep to defer to the leadership of their elders and suffer themselves to be led by these God-ordained men.  You have your freedom, and you may use it as you will.  But you will be held to account for \textit{how} you use your freedom.  

As you exercise your freedom in a church election, ask yourself:
\begin{quote}
\begin{enumerate}
    \item How would God want me to vote on this matter?
    \item Am I obeying my leaders and submitting to them in how I am casting my vote?
    \item Have I submitted to the leadership of our elders in doing my due diligence (prayer, study, research, examination) before coming to cast my vote?
    \item Do I have a clear \textit{biblical} case to vote the way I'm intending?
\end{enumerate}
\end{quote}

\noindent
Alexander Strauch reminds us:
\begin{quote}
Conflict between leaders and the led can at times become severe.  Ultimately, however, God uses these discordant situations to expose our pride, selfishness, and lack of love.  Believers who love their leaders will have greater understanding and tolerance for their mistakes or the decisions they don't understand or agree with.
\end{quote}

\noindent
He continues:
\begin{quote}
In love, believers will view difficult situations in the best possible light.  In love, believers will be less critical and more responsive to the elders' instruction and admonition.\footcite{strauch2024esteem}
\end{quote}

\noindent
Finally, regarding the issue of voting for or recognizing elders, Strauch has this to say:
\begin{quote}
... perpetuation of the eldership is implied in the elders' role as community leaders, shepherds, stewards, and overseers.  Perpetuating the eldership is a major aspect of church leadership responsibility.  It is absolutely vital to the life of the church and the eldership that the elders recognize the desire of others to shepherd the flock.  It is all too easy to allow personal feelings or fear to prejudice one's judgment toward aspiring elders.  \textit{One must have biblical reasons to keep a man from church oversight}.  Yet many qualified men are refused their rightful place as elders because they are a threat to the status quo.  Indeed good men are often a threat to the status quo.  Far too often, church politics rule in these matters instead of God's Word, prayer, and the leading of God's Spirit.\footcite[p. 75]{strauch2000recovering}
\end{quote}

\noindent
The operative part of this quote is italicized in the original, ``\textit{one must have biblical reasons to keep a man from church oversight}''.  Not a vague feeling, not a personal preference.  \textit{Biblical reasons}.  As we will discuss next, you must exercise the \textit{Duty of Diligence} with regard to all matters coming before the congregation, meaning that you must search out the matter and inform yourself about what's at stake.  In the course of exercising your due diligence, you may discover adverse information about the matter, which you have a \textit{Duty to Declare} to your elders.

\ifdefined\embedversion\newpage\fi
\section{Duty of Diligence}
Diligence means having a persevering determination to perform a task.  It is a spirit of conscientiousness in paying proper attention to something, or giving the degree of care required in a given situation.  God's people are called to be diligent.  Solomon uses the Hebrew word transliterated into our alphabet as \textit{charuwts} throughout Proverbs, and this is translated as "diligent" in many of our English translations.  This word "gathers several concrete and figurative ideas into a single word family: the cutting edge of a threshing sledge, the priceless sheen of refined gold, the unwavering resolve of a royal decree, and the focused energy of diligent labor."\footcite{biblehub2742}  

Let's consider some quotes from Solomon on the matter:
\begin{quote}
Poor is he who works with a negligent hand, But the hand of the \textbf{diligent} makes rich.
\\ \hspace*{\fill} --Proverbs 10:4 (\textsc{nasb 77})

The hand of the \textbf{diligent} will rule, But the slack hand will be put to forced labor.
\\ \hspace*{\fill} --Proverbs 12:24 (\textsc{nasb 77})

A slothful man does not roast his prey, But the precious possession of a man is \textbf{diligence}.
\\ \hspace*{\fill} --Proverbs 12:27 (\textsc{nasb 77})

The soul of the sluggard craves and gets nothing, But the soul of the \textbf{diligent} is made fat.\footnote{I promise you this was considered a good thing in this context!}
\\ \hspace*{\fill} --Proverbs 13:4 (\textsc{nasb 77})

The plans of the \textbf{diligent} lead surely to advantage, But everyone who is hasty comes surely to poverty.
\\ \hspace*{\fill} --Proverbs 21:5 (\textsc{nasb 77})
\end{quote}

Solomon contrasts diligence with \textit{negligence}.  He contrasts the success of the diligent with the failure of the \textit{slack hand}.  Diligence is portrayed as an antonym of \textit{laziness}.  Interestingly, also, diligence is held to be far from \textit{everyone who is hasty}, which is interesting, because it seems to indicate that haste, while it may seem to be out of the wheelhouse of the physically lazy man - the \textit{mentally} lazy man might be forced to be hasty because he waits until the last minute to prepare.

God's people are called in Deuteronomy to "\textbf{diligently} keep the commandments of the \textsc{Lord} your God, and His testimonies and His statutes which He has commanded you."\footnote{Deuteronomy 6:17 (\textsc{nasb 77})}  Deuteronomy also advises us to "give heed to yourself and keep your soul \textbf{diligently}, lest you forget the things which your eyes have seen ... but make them known to your sons and your grandsons."\footnote{Deuteronomy 4:9 (\textsc{nasb 77})}

Finally, the writer of Hebrews says:
\begin{quote}
And we desire that each one of you show the same \textbf{diligence} so as to realize the full assurance of hope until the end, that you may not be sluggish, but imitators of those who through faith and patience inherit the promises.
\\ \hspace*{\fill} --Hebrews 6:11-12 (\textsc{nasb 77})
\end{quote}

As Christians, we are called to be discerning, and we are called to exercise judgment.  When we come to a business meeting, we are conducting the Lord's business for His local body, and we should come prepared.  We should be diligent in our preparation.  We are called to wisdom, and we are called to ask God if we lack wisdom.  There are many Bible verses illustrating this reality.\footnote{I Thessalonians 5:21; I Timothy 3:10; Hebrews 13:7; I Corinthians 6:1-8 (in the context of conflicts); Romans 12:2; I John 4:1; James 1:5; Hebrews 5:14}

What if we don't have enough information about a matter?  We can ask questions, we can research it, we can study, we can pray.  In other words, we can exercise \textit{diligence} with regard to the matter of informing ourselves about church matters.  We have heard in recent government election cycles about so-called "low information voters".  Don't allow yourself to be viewed as an example of this phenomenon.  Inform yourself.  Exercise the diligence that the matter is due.  And if you do discover information adverse to the way you're being directed to vote, you have a \textit{Duty to Declare} that information to the elders.

\ifdefined\embedversion\newpage\fi
\section{Duty to Declare}
Christian church members have a \textit{Duty to Declare} any material information they may know about any matters being brought before the church.  When information is sensitive or scandalous, such declaration should happen only to the elders.  Sometimes, however, it should happen to the body at large.  This duty is \textit{especially} important if a member has information about a matter which is going to lead them to vote \textit{against} the leadership of the elders on a matter.

The constitution of Capitol Hill Baptist Church in Washington, DC, where Mark Dever of 9Marks serves, says that the congregation should not vote against recognizing an elder candidate \textit{unless} they have spoken about their concern with their elders.\footcite[\\The full relevant quote from that page is, "In fact, our church constitution says that we can only vote against recognizing a prospective elder if we've spoken with an elder ahead of time about our concern.  \textbf{The reason for this is simple: if there's some concern you have about that individual's ability to lead this congregation that's significant enough that you would withhold your vote, it might be a good reason for the elders to withdraw that person's nomination}."]{chbc2016governance} The actual text of their constitution on this matter is found in Article 6, Section 2, Paragraph 2:
\begin{quote}
The elders should seek recommendations and involvement from the general membership in the nomination process. \textbf{Any member with reason to believe that a nominated candidate is unqualified for an office \textit{should} express such concern to the elders}. Members intending to speak in opposition to a candidate should express their objection to the elders \textit{as far in advance as possible before the relevant church members' meeting}.\footcite[\\The emphasis here is mine.]{chbc2021constitution}
\end{quote}

Elders aren't perfect.  As I've noted, they have blind spots.  They don't know everything.  If you have information about an elder candidate, about somebody becoming a member, about entering into a building project, and you don't provide that information to your elders, you are failing in a key responsibility.  You are possibly opening up your elders to ridicule.  You are possibly paving the way to a very disastrous and church-splitting vote.  

Consider the golden rule.\footnote{Therefore, however you want people to treat you, so treat them, for this is the Law and the Prophets. \\ \hspace*{\fill} --Matthew 7:12 (\textsc{nasb 77})}  If \textit{you} were an elder and had commended Joe Blow to the congregation for the office of elder, but you had no idea about Mr. Blow's somewhat salacious social media presence, and the elder vote failed because about half of your congregation \textit{did} know but failed to report it to you, how would you feel?  Would you feel like the work was bringing you joy or grief?\footnote{Hebrews 13:17, again.} What kind of a testimony would such a show of disunity display to the outside world?

\ifdefined\embedversion\newpage\fi
\section{Duty to Decline}
Christian church members have a \textit{Duty to Decline} or abstain from voting if they cannot in good conscience vote in the way their elders are leading \textit{if} they have contrary information but have not brought it forward.  In other words, if they have not exercised their \textit{Duty to Declare}, they \textit{should} exercise their \textit{Duty to Decline}.

This follows logically from two other duties - the \textit{Duty to Defer} and the \textit{Duty to Declare}. If you cannot defer to your elders but you have not brought forward information injurious to the direction in which the elders are leading, then you: 1) should not do harm to your conscience by voting \textit{with} your elders, and 2) should not do harm to your relationship with your elders by voting \textit{against} your elders without having exercised your \textit{Duty to Declare}.  The safest option in this case, I believe, is to decline to vote, or to abstain.\footcite{dever2024vote}  This is the most inferential duty mentioned so far, so I hesitate to say that you \textit{must} exercise this duty, although I would say you \textit{must} prayerfully consider this course of action in such a situation in order to preserve your conscience.  Abstention is a valid way to resolve the tension between voting your conviction and following your elders' lead.

\ifdefined\embedversion\newpage\fi
\section{Duty to Depart}
Christian church members have a \textit{Duty to Depart} a body if they find themselves in regular, fundamental disagreement with their elders.  If such disagreements cannot be resolved, you are putting yourself, your family, and your elders in a difficult situation which is not likely to be pleasant for anyone involved.

In the \textit{Pastors Talk} episode cited previously, there is an exchange between Jonathan Leeman and Mark Dever:
\begin{quote}
\textbf{Jonathan Leeman:}
So Mark, for you same answer [question] on this one. "Look, I've been following this guy, I've been following him on Twitter, he's a jerk and I think he might be deconstructing. It's not really clear, but now the elders want me to vote for him." You're again saying --

\textbf{Mark Dever:}
Either your elders are whacked or they just don't know something you know.

\textbf{Jonathan Leeman:}
Yeah.

\textbf{Mark Dever:}
Yeah. If they're whacked, you need to leave. You need to leave the church. And \textit{if they don't know, it's your fault for not telling them}.\footcite[emphasis mine]{dever2024vote}
\end{quote}

Of course it may not be the elders who are \textit{whacked}, as Dever puts it.  Titus 3:10-11 has words for the "factious man", that he should be rejected, and that he is in sin.\footnote{Reject a factious man after a first and second warning, knowing that such a man is perverted and is sinning, being self-condemned. \\ \hspace*{\fill} --Titus 3:10-11 (\textsc{nasb 77})}  We need to ask ourselves when we are in conflict (especially repeatedly) if \textit{we} are the problem.

As Christians, we are called to peace.\footnote{If possible, so far as it depends on you, be at peace with all men.\\ \hspace*{\fill} --Romans 12:18 (\textsc{nasb 77})}  We are called to be unified with our brethren.\footnote{To sum up, let all be harmonious, sympathetic, brotherly, kindhearted, and humble in spirit; not returning evil for evil, or insult for insult, but giving a blessing instead; for you were called for the very purpose that you might inherit a blessing.\\ \hspace*{\fill} --I Peter 3:8-9 (\textsc{nasb 77})}  If you find yourself in a situation where you are unable to live in peace and unity with your local body, and especially with your elders, you need to prayerfully consider leaving.  It may be you, it may be them, it may be some combination.  But ultimately, peace and unity must be preserved.

Do not make such a decision flippantly or for fleshly reasons.  Do not make an abrupt departure in which you simply disappear or "ghost" your fellow church members.  Let people know \textit{why} you're leaving without betraying any confidences or making the division any worse than it needs to be.  \textit{How} you leave reflects on you as much or more as \textit{why} you leave.

Also, do not leave without a destination in mind.  Christians were not designed to be detached from a local body of Christ.  Your faith will shrivel without Christian fellowship.  

\ifdefined\embedversion\newpage\fi
\section{Duty to Disencumber}
Christian church elders have a \textit{Duty to Disencumber} the congregation so that they can provide meaningful feedback through the voting process.\footcite{strauch2024elderrelation}  They should be even-handed and fair in their presentation of issues up for consideration so as to get as authentic a result as possible, or else congregational input means nothing.  The congregation should not be encumbered, or burdened, by pressure from the elders.

This duty exists in a healthy tension with the \textit{Duty to Defer}.\footnote{This tension is properly resolved through mutual love -- the elders' love for the congregation leads them to avoid being overbearing, and the congregation's love for the elders leads them to defer to their God-ordained leaders.  This is all orchestrated by the whole congregation's mutual love for God and His ways.}  The congregation \textit{must} defer to its elders, but the elders \textit{must} deal with the congregation in good faith, earnestly teaching and convincing rather than manipulating people or circumstances.  This issue is addressed in the \textit{Pastors Talk} episode quoted from before:
\begin{quote}
\textbf{Jonathan Leeman:}
It needs to be a real decision.

\textbf{Jamie Dunlop:}
It needs to be a real decision. And elders, especially elders who are trusted, need to be very careful in how they steward that trust.

\textbf{Mark Dever:}
This is why with church discipline cases, even with the clearest ones, we try to give a level of information to the congregation as a whole so that they would have enough information to vote against our recommendation should they want to do that.

\textbf{Jamie Dunlop:}
We don't hide the contrary information.

\textbf{Jonathan Leeman:}
Because that robs the decision of integrity for them.

\textbf{Jamie Dunlop:}
Yeah. And we don't tell them, "Now, this really is a Hebrews 13:17 situation. You better vote with us." Even if I might feel that way based on this podcast.\footcite{dever2024vote}
\end{quote}

Just as a man would not remind his wife of Ephesians 5:22-24 in the midst of a disagreement,\footnote{Wives, be subject to your own husbands, as to the Lord. For the husband is the head of the wife, as Christ also is the head of the church, He Himself being the Savior of the body. But as the church is subject to Christ, so also the wives ought to be to their husbands in everything. \\ \hspace*{\fill} --Ephesians 5:22-24 (\textsc{nasb 77})} so elders shouldn't browbeat the sheep with admonitions from Hebrews 13:17 or elsewhere in the midst of a business meeting.  It wouldn't work out well for the first situation,\footnote{Ask me how I know!} and it won't work out well for the second.  It will result in an erosion of trust, with predictable consequences.

Additionally, if you tell your congregation that their vote matters, then it's up to you as elders to keep that promise and make sure that their vote really does matter.  Resist the urge to manipulate a response.  

That said, every decision does not need to be up for a vote, and removing the option to cast a vote for certain matters might be one way to disencumber the congregation. For example, the ministry direction of the church should be chosen by the elders, after much prayer and study, by the leading of God.  \textit{However}, even though the elders have a right to govern church activities, \textit{the degree to which they do this without getting buy-in from the congregation is the degree to which they risk alienating them}.\footnote{Similarly, although members have a right to vote how they please, voting against the elders without bringing concerns to them beforehand erodes trust.}

According to the Second London Baptist Confession of Faith, elders and deacons are to be chosen by the "common suffrage of the church it self".\footcite[Chapter 26, Paragraph 9]{london1689}  While I find clear evidence in the New Testament of voting for the office of deacon,\footnote{Acts 6:1-6} I cannot find clear evidence of a church casting votes for the office of elder.  Rather, the New Testament pattern seems to be godly men appointing godly men.  

That said, I can appreciate the \textit{prudence} of elders governing by the consent of the governed.  Again the principle applies: \textit{the degree to which elders chart ministry direction without getting buy-in from the congregation is the degree to which they risk alienating them}.  But even when a formal vote is conducted for the office of elder, God's people should not understand such an election as a purely competitive election in a democratic sense, but rather as an opportunity to recognize or affirm the qualifications of men which have already been vetted and chosen by the elders to lead.

For matters which \textit{are} to be voted on, both Scripture and the Second London Baptist Confession of Faith are silent as to a method.  The most common types of voting in Baptist churches seem to be raising of hands\footnote{If you want to see Baptists raising hands, you have to go to a business meeting.} and, in more modern times, secret ballots.  There are ditches on both sides of the road.

Raising of hands in a public assembly to vote on a church matter can result in people voting for social reasons instead of spiritual or practical ones.  Public or open voting can also theoretically result in retaliation if any members are in a position to retaliate against other members based on their vote.  In my opinion, secret ballots, while solving some of the problems of public voting, have the potential to cause more problems than they solve.  First, many will get the impression that this is "just like" a governmental election, that their vote is their personal choice, and that they are accountable only to themselves.  Second, it can stifle any sense of openness within a church body, as people are loathe to talk about votes cast by way of a "secret"\footnote{Additionally, in small churches, it is not very realistic to believe that secret ballots will really \textit{remain} secret.  For someone who knows the congregation, it's a trivial process of elimination to determine exactly how each member voted.} ballot.  I believe these concerns of decreased accountability and openness are of greater import than concerns about "social voting" or retaliation.  If people are engaging in such un-Christian activities as a result of a public ballot, at least these actions are done in the open, which yields greater potential for accountability measures to be carried out.

\ifdefined\embedversion\newpage\fi
\section{Conclusion}
\newthought{Ultimately, each church must decide} for itself on what matters it will vote and how it will conduct such voting, using Scripture as a guide.  Elders are responsible for making sure the process is fair, consistent, and honest, but members are responsible for deferring to the leadership of the elders.  Plans to vote adversely to the elders must be communicated, and if they are not communicated, the member should consider abstention as a way to preserve conscience with regard to their convictions as well as with regard to their relationship with their church leaders.  Persistent disagreements with elders should be a cause for concern, and \textit{may} be a cause to depart, but they don't have to be.  

Special care is required to realize the magnitude of the tension which can exist between the \textit{Duty to Defer} and \textit{Duty to Disencumber}.  Members may believe they have a duty to trust and defer to their elders, but yet believe (rightly) they have a duty to act faithfully and honor God instead of simply outsourcing their judgment to the elders.  Similarly, elders may think that a particular course of action is best for the church, but also believe that they may not coerce consciences.  The healthiest solution to this conundrum is for members to recognize that they need a positive reason \textit{not} to follow their elders, and for elders to realize that they need to persuade the congregation rather than pressuring them.  The more strongly deference is emphasized, the more congregational responsibility may be undermined.  Similarly, the more strongly congregational independence is emphasized, the more the concept of elder-led shepherding can suffer.  The right balance can only be arrived at through prayer and mutual affection within the body.

Churches are \textit{not} democracies, no matter how congregational their polity might be.  The Kingdom of God is a monarchy.  Christ is King.  Let us never get the incorrect idea that church elections are solely up to individual preference and that church members are not liable to God and their fellow church members for doing their duties when voting on matters before the church.  We must realize that Christians will be held to account for how they vote on church matters, whether the vote is secret or not.  Elders will be held to account to an even higher degree for how they lead their flocks.  We must vote, and lead, in a way that glorifies our Father in Heaven. 
\end{document}
